\documentclass[11pt, a4paper]{article}

% ---- Packages ----
\usepackage[utf8]{inputenc}
\usepackage[T1]{fontenc}
\usepackage{amsmath, amssymb, amsthm}
\usepackage{graphicx}
\usepackage{booktabs}
\usepackage{hyperref}
\usepackage[margin=1in]{geometry}
\usepackage{natbib}
\usepackage{enumitem}
\usepackage{listings}
\usepackage{xcolor}
\usepackage{authblk}
\usepackage{abstract}

% ---- Formatting ----
\hypersetup{
    colorlinks=true,
    linkcolor=blue,
    citecolor=blue,
    urlcolor=blue
}

\lstset{
    basicstyle=\ttfamily\small,
    breaklines=true,
    frame=single,
    backgroundcolor=\color{gray!5},
    columns=fullflexible,
    keepspaces=true
}

\renewcommand{\abstractnamefont}{\normalfont\bfseries}
\renewcommand{\abstracttextfont}{\normalfont\small}

\newtheorem{definition}{Definition}
\newtheorem{theorem}{Theorem}
\newtheorem{corollary}{Corollary}

% ---- Metadata ----
\title{\textbf{From MEV to GEV: Generalized Extractable Value Resistance\\Through Cooperative Mechanism Design}}

\author[1]{William Glynn}
\affil[1]{VibeSwap --- Independent Research}
\date{March 2026}

% ============================================================
\begin{document}
\maketitle

% ---- Abstract ----
\begin{abstract}
Maximal Extractable Value (MEV) has dominated DeFi security discourse since its formalization by Daian et al.\ (2019). Solutions---private mempools, MEV auctions, encrypted order flow---treat MEV as an isolated pathology. We argue MEV is one instance of a broader class we term \textbf{Generalized Extractable Value (GEV)}: any value that an intermediary can capture from a system's participants by exploiting structural asymmetries in information, governance, timing, or capital access. We identify seven distinct GEV vectors present across major DeFi protocols and present VSOS (VibeSwap Operating System), a financial operating system where each layer structurally eliminates its corresponding extraction vector through Shapley-fair attribution, physics-based parameter control, and non-rent-seeking token economics. The central claim is architectural: MEV-resistance is a feature; GEV-resistance is a design principle that must be applied at every layer of a financial system, not bolted onto one.
\end{abstract}

\noindent\textbf{Keywords:} MEV, extractable value, mechanism design, Shapley values, cooperative game theory, DeFi, rent-seeking

% ============================================================
\section{Introduction}

Every major DeFi protocol solves a real problem. Chainlink solved oracle reliability. AAVE solved permissionless lending. Curve solved low-slippage stablecoin trading. Synthetix solved synthetic asset issuance. Then each one put a tollbooth in front of the solution.

LINK holders extract rent from oracle queries without providing data. AAVE governance token holders extract protocol revenue without providing liquidity. CRV's vote-escrow mechanism created a permanent governance aristocracy monetized through bribe markets. SNX stakers bear socialized debt risk without Shapley-proportional attribution of their marginal contribution.

These are not bugs. They are features---features designed to capture value for intermediaries at the expense of the users who generate it. MEV is merely the most visible instance of this pattern because it operates on millisecond timescales and leaves on-chain evidence. But the extraction is the same whether it takes milliseconds (sandwich attacks), months (VC lockup dumps), or is continuous (token rent on every transaction).

We propose that the correct unit of analysis is not MEV but \textbf{GEV}---the total extractable value across all structural asymmetries in a protocol. A system is GEV-resistant when no participant receives more than their marginal contribution warrants, as measured by Shapley values from cooperative game theory \citep{shapley1953}.

% ============================================================
\section{Taxonomy of Generalized Extractable Value}

We identify seven distinct GEV vectors, each exploiting a different structural asymmetry.

\subsection{Transaction Ordering EV (MEV)}

\textbf{Asymmetry:} Information about pending transactions in the mempool. Validators, builders, and searchers reorder, insert, or censor transactions to capture value. Flashbots reports cumulative MEV extraction exceeding \$680M on Ethereum mainnet through 2024 \citep{flashbots2024}.

\subsection{Governance Extractable Value (GoEV)}

\textbf{Asymmetry:} Concentration of governance tokens enables minority control over extractable parameters. Bribe markets (Convex, Votium) convert public goods---emission direction---into private rent. MakerDAO's stability fee is set by MKR holders who benefit from higher fees at the expense of borrowers.

\subsection{Token Rent-Seeking EV (TrEV)}

\textbf{Asymmetry:} Protocol-mandated token intermediation in value flows that don't require a token. LINK must be held by node operators not because the oracle mechanism requires it, but because the tokenomics model imposes it as economic rent. GRT curators stake tokens to signal subgraph quality---the staking is speculative positioning, not a quality signal.

\subsection{Capital Formation EV (CfEV)}

\textbf{Asymmetry:} Asymmetric access to token supply at formation. The median VC-backed token in 2024 experienced a 72\% price decline within 12 months of lockup expiry \citep{messari2024}. The value extracted is the delta between VC entry price and average retail purchase price, multiplied by supply sold.

\subsection{Oracle Extractable Value (OrEV)}

\textbf{Asymmetry:} Control over the data pipeline between off-chain reality and on-chain state. Oracle providers impose rent via mandatory token intermediation. Oracle latency creates arbitrage: anyone who sees the off-chain price before the on-chain update can trade the stale price.

\subsection{Platform Extractable Value (PlEV)}

\textbf{Asymmetry:} Platform control over user relationships, data, and access. Centralized frontends, proprietary APIs, and data silos enable operators to extract value from user activity through fee capture, data harvesting, and access gating.

\subsection{Liquidation Extractable Value (LqEV)}

\textbf{Asymmetry:} Priority access to liquidation transactions. On Hyperliquid and dYdX, the sequencer has privileged access to liquidation ordering. Liquidation bots compete through MEV infrastructure, extracting value from borrowers who receive worse liquidation prices than market clearing would produce.

% ============================================================
\section{The GEV Framework}

\subsection{Formal Definition}

Let $P$ be a protocol with participant set $N = \{1, 2, \ldots, n\}$. For each participant $i$, define $v_i$ as the value $i$ generates, $\phi_i$ as $i$'s Shapley value, and $r_i$ as the value $i$ actually receives.

\begin{definition}[Generalized Extractable Value]
The GEV of protocol $P$ is:
\begin{equation}
    GEV(P) = \sum_{i \in N} \max(0,\; r_i - \phi_i)
    \label{eq:gev}
\end{equation}
A protocol is \textbf{GEV-resistant} iff $GEV(P) = 0$.
\end{definition}

\subsection{Properties}

\textbf{Comprehensiveness.} MEV captures only transaction-ordering extraction. GEV captures all structural extraction.

\textbf{Shapley grounding.} GEV is measured against the Shapley value---the unique allocation satisfying efficiency, symmetry, linearity, and the null player property \citep{shapley1953}.

\textbf{Structural, not behavioral.} GEV measures extraction potential embedded in protocol design, not individual bad behavior.

\begin{corollary}[Conservation of Extraction]
Eliminating GEV at one layer without addressing all layers relocates extraction. Partial GEV-resistance is not GEV-resistance.
\end{corollary}

This is the fundamental insight: \textbf{extraction is conserved across layers}. MEV-resistant protocols with rent-seeking governance tokens have not reduced GEV---they have relocated it.

% ============================================================
\section{VSOS: Layer-by-Layer GEV Elimination}

VSOS (VibeSwap Operating System) is designed for $GEV = 0$ across all layers. Each vector is addressed by a specific mechanism:

\begin{table}[h]
\centering
\caption{GEV vectors and VSOS defenses}
\label{tab:defenses}
\begin{tabular}{@{}lll@{}}
\toprule
\textbf{GEV Vector} & \textbf{VSOS Defense} & \textbf{Mechanism} \\
\midrule
MEV & Commit-reveal batch auctions & Sealed-bid, uniform clearing price \\
GoEV & Ungovernance + PID control & Voting decay, physics-based params \\
TrEV & Shapley distribution & No tollbooth tokens, pay-per-use \\
CfEV & Fair launch + ContributionDAG & Retroactive Shapley, no VC rounds \\
OrEV & Quality-weighted oracle rewards & Accuracy scoring, no LINK rent \\
PlEV & Fork escape + mutualist absorption & Zero-cost exit, cooperative model \\
LqEV & Batch-settled perpetuals & Same commit-reveal for liquidations \\
\bottomrule
\end{tabular}
\end{table}

\subsection{MEV: Commit-Reveal Batch Auctions}

Orders are submitted as hashed commitments during an 8-second commit phase, revealed during a 2-second reveal phase, then settled at a uniform clearing price after a Fisher-Yates shuffle using XOR'd participant secrets. No participant can observe others' orders before committing (information symmetry). The uniform clearing price eliminates intra-batch price discrimination.

\subsection{GoEV: Ungovernance + PID Auto-Tuning}

The Ungovernance Time Bomb decays all voting power over time. Parameters that typically require governance votes---stability fees, emission rates, interest rate curves---are controlled by PID controllers that auto-adjust based on on-chain signals. If parameters are set by physics, there are no parameters for governance to capture.

\subsection{TrEV: Shapley Distribution}

Protocol revenue is distributed via Shapley values computed on-chain. No tokens must be held to access services. The pairwise Shapley verification formula enables $O(1)$ on-chain checking:

\begin{equation}
    |\phi_i \cdot w_j - \phi_j \cdot w_i| \leq \varepsilon
    \label{eq:pairwise}
\end{equation}

Any participant can verify their allocation is fair without trusting the protocol.

\subsection{CfEV: Fair Launch + ContributionDAG}

No VC rounds. No insider allocation. Distribution is via contribution mining tracked in an on-chain ContributionDAG. The Cave Theorem states that foundational contributions earn more by the mathematics of Shapley values (they appear in more coalitions), not by timestamp advantage.

\subsection{OrEV: Quality-Weighted Oracle Rewards}

VibeOracleRouter aggregates first-party feeds (API3 pattern), aggregated feeds (Chainlink pattern), and low-latency feeds (Pyth pattern). Providers are rewarded by accuracy contribution via Shapley attribution. Accuracy scores decay over time to prevent lock-in.

\subsection{PlEV: Fork Escape + Mutualist Absorption}

The Fractal Fork Network ensures any layer can be forked at zero cost. This implements Hirschman's exit-voice framework \citep{hirschman1970} as a protocol primitive. If exit is costless, the only sustainable strategy is to provide value at marginal cost---which is the $GEV = 0$ condition.

\subsection{LqEV: Batch-Settled Perpetuals}

VibePerpEngine settles liquidations through commit-reveal batch auctions. No liquidator has ordering advantage. The liquidation discount is set by market clearing, not by bot speed.

% ============================================================
\section{The Composability Constraint}

GEV-resistance is not compositional by default. Two GEV-resistant modules can create GEV when composed if their interface permits extraction. VSOS addresses this through three composability constraints:

\begin{enumerate}[leftmargin=*]
    \item \textbf{Unified attribution}: All modules read \texttt{ShapleyDistributor} for reward distribution.
    \item \textbf{Unified credit}: All modules write to \texttt{ContributionDAG} for attribution tracking.
    \item \textbf{Unified safety}: All modules use \texttt{CircuitBreaker} for volume, price, and withdrawal bounds.
\end{enumerate}

If every value flow passes through one attribution mechanism, every credit is recorded in one graph, and every risk is bounded by one breaker, then no cross-module extraction is possible that isn't caught by one of the three.

% ============================================================
\section{Comparison with Existing Approaches}

\begin{table}[h]
\centering
\caption{GEV reduction by approach}
\label{tab:comparison}
\begin{tabular}{@{}lcc@{}}
\toprule
\textbf{Approach} & \textbf{Vectors Addressed} & \textbf{GEV Reduction} \\
\midrule
Flashbots / MEV-Share & MEV only & $\sim$14\% \\
Encrypted mempools (Shutter) & MEV only & $\sim$14\% \\
Batch auctions (CowSwap) & MEV only & $\sim$14\% \\
Conviction voting & GoEV (partial) & $<$14\% \\
Fair launch (Balancer LBP) & CfEV (partial) & $<$14\% \\
\textbf{VSOS} & \textbf{All 7 vectors} & \textbf{100\%} \\
\bottomrule
\end{tabular}
\end{table}

The differentiator is not that VSOS has better solutions for any single vector. The differentiator is that VSOS applies GEV-resistance as an architectural constraint across every layer, ensuring extraction cannot relocate.

% ============================================================
\section{Limitations and Open Questions}

\subsection{Shapley Computation}
Exact Shapley values require evaluating $2^n$ coalitions. VSOS uses off-chain computation with on-chain $O(1)$ pairwise verification (Eq.~\ref{eq:pairwise}). The security of this approach against adversarial Shapley reports remains an open challenge.

\subsection{PID Controller Robustness}
PID controllers can be manipulated through their input signals. An attacker who influences the on-chain observables (utilization, peg deviation) can indirectly control parameters. TWAP validation (max 5\% deviation) and circuit breakers mitigate this, but formal robustness guarantees against adversarial PID manipulation remain open.

\subsection{Fork Credibility}
The fork escape argument assumes rational actors who fork when extraction exceeds forking costs. If network effects make forking costly despite being technically free, some PlEV may persist---the same limitation that constrains competition in traditional markets.

\subsection{Cross-Chain GEV}
Cross-chain message latency via LayerZero V2 creates a new vector: actors who observe state on one chain can act on another before the message arrives. BatchProver's cross-chain proof verification reduces but may not eliminate this vector.

% ============================================================
\section{Conclusion}

MEV is not the disease. It is a symptom. The disease is structural extraction---value captured by intermediaries through asymmetries in information, governance, timing, capital access, data control, and liquidation priority.

We have formalized this as Generalized Extractable Value and identified seven distinct vectors. We have shown that partial solutions merely relocate extraction to unaddressed vectors. Extraction is conserved across layers.

VSOS demonstrates that GEV-resistance is achievable as an architectural property. The mechanisms are not novel individually---batch auctions, Shapley values, PID controllers, and fair launches are established primitives. The contribution is compositional: applying all of them simultaneously, at every layer, with three composability constraints that ensure GEV-resistance is preserved under composition.

The result is a financial operating system where the answer to ``who extracts?'' is ``nobody''---not because extraction is prohibited by governance (which can be captured), but because extraction is impossible by construction. Policy becomes physics.

\begin{quote}
\emph{Every protocol on the VSOS absorption list solved a real problem. Then they put a tollbooth in front of the solution. We remove the tollbooth and replace it with math.}
\end{quote}

% ============================================================
\bibliographystyle{plainnat}
\begin{thebibliography}{10}

\bibitem[Daian et~al.(2019)]{daian2019}
Daian, P., Goldfeder, S., Kell, T., et~al. (2019).
\newblock Flash Boys 2.0: Frontrunning, transaction reordering, and consensus instability in decentralized exchanges.
\newblock \emph{IEEE Symposium on Security and Privacy}, 2020.

\bibitem[Shapley(1953)]{shapley1953}
Shapley, L.~S. (1953).
\newblock A value for n-person games.
\newblock \emph{Contributions to the Theory of Games II}, Annals of Mathematics Studies 28, pp.~307--317.

\bibitem[Hirschman(1970)]{hirschman1970}
Hirschman, A.~O. (1970).
\newblock \emph{Exit, Voice, and Loyalty: Responses to Decline in Firms, Organizations, and States.}
\newblock Harvard University Press.

\bibitem[Roughgarden(2021)]{roughgarden2021}
Roughgarden, T. (2021).
\newblock Transaction fee mechanism design.
\newblock \emph{ACM Conference on Economics and Computation}.

\bibitem[{Flashbots}(2024)]{flashbots2024}
{Flashbots} (2024).
\newblock MEV-Explore: Cumulative extracted MEV.
\newblock \url{https://explore.flashbots.net}.

\bibitem[{Messari}(2024)]{messari2024}
{Messari} (2024).
\newblock VC token unlock impact analysis.
\newblock Messari Research Report.

\bibitem[Buterin(2021)]{buterin2021}
Buterin, V. (2021).
\newblock Moving beyond coin voting governance.
\newblock \url{https://vitalik.eth.limo}.

\bibitem[Glynn(2026a)]{glynn2026vsos}
Glynn, W. (2026).
\newblock VSOS protocol absorption: 28 protocols into 9 modular layers.
\newblock VibeSwap Documentation.

\bibitem[Glynn(2026b)]{glynn2026symbolic}
Glynn, W. (2026).
\newblock Symbolic compression in human-{AI} knowledge systems.
\newblock arXiv preprint.

\bibitem[Glynn(2018)]{glynn2018}
Glynn, W. (2018).
\newblock Wallet security axioms: Seven principles for cryptocurrency key management.
\newblock Independent Publication.

\end{thebibliography}

\end{document}
